% TOP_PROBLEM: {"id": 3125, "title": "Do any three longest paths in a connected graph have a vertex in common?", "category_id": 3, "category": {"id": 3, "name": "graph_theory", "display_name": "Graph Theory", "description": "Problems involving graphs, networks, and their properties.", "slug": "graph-theory", "order_index": 3, "created_at": "2026-07-31T15:26:25.671Z"}, "status": "open", "problem_number": "OPG-46538", "set_id": 12}
% FIRST_POSTED: 2026-10-01
% ENTRY_KIND: statement_only
% UnsolvedMath ID 3125; source catalogue code OPG-46538
% !TeX program = lualatex
% Definition and sources only; this file is not an LLM solution attempt.
\documentclass[11pt,a4paper]{article}
\usepackage[margin=25mm]{geometry}
\usepackage{fontspec}
\setmainfont{lmroman10-regular.otf}[BoldFont=lmroman10-bold.otf,
 ItalicFont=lmroman10-italic.otf,BoldItalicFont=lmroman10-bolditalic.otf]
\usepackage{amsmath,amssymb,mathtools}
\usepackage{unicode-math}
\setmathfont{latinmodern-math.otf}
\AtBeginDocument{\let\setminus\smallsetminus}
\usepackage{xurl,hyperref}
\hypersetup{colorlinks=true,urlcolor=blue}
\setcounter{secnumdepth}{0}
\setlength{\parindent}{0pt}
\setlength{\parskip}{5pt}
\setlength{\emergencystretch}{3em}
\begin{document}
\section*{Do any three longest paths in a connected graph have a vertex in common?}\label{3125}
{\small UnsolvedMath ID 3125 · OPG-46538 · Graph Theory · OpenGarden}
\subsection{Definitions and mathematical statement}
Let $G=(V,E)$ be a finite nonempty connected simple undirected graph.
A path is a sequence $v_0,\ldots,v_t$ of distinct vertices with an edge
between each consecutive pair; its length is $t$, its number of edges.
Define $L(G)$ to be the largest length of a path in $G$.
A longest path is any path attaining $L(G)$.

The conjecture states that for every such $G$ and every three longest
paths $P_1,P_2,P_3$,
\[
 V(P_1)\cap V(P_2)\cap V(P_3)\ne\varnothing.
\]
The common vertex may depend on the chosen triple. It need not be an
endpoint of any path, and no common edge is required. Allowing repetition
among the three paths does not strengthen the distinct-path case.
For a graph with one vertex, the length-zero path uses that vertex.

The assertion is about a triple's simultaneous intersection, which does
not follow merely from each pair having a nonempty intersection. It also
does not assert that a single vertex belongs to every longest path of
the graph: different triples can have different witnesses. Paths are
not required to be induced; edges between nonconsecutive path vertices
are permitted in $G$.

``Longest'' refers to maximum length among all paths in the entire graph,
not just to inclusion-maximal paths or to shortest paths between their
ends. A negative resolution would consist of a connected graph and
three globally longest paths whose three vertex sets have empty
intersection.
\subsection{Short English statement}
In every connected graph, must any selected three longest paths share at least one vertex?
\subsection{Sources}
\begin{itemize}
\item[S1] ULAM.AI and the UnsolvedMath Contributors, supplied problems.json, record 3125 (OPG-46538), OpenGarden collection. Catalogue identity and source transcription; CC BY 4.0.
\url{https://huggingface.co/datasets/ulamai/UnsolvedMath}.
\item[S2] Open Problem Garden, Do any three longest paths in a connected graph have a vertex in common?. Original problem and discussion; the mathematical formulation below is expanded from the supplied source record.
\url{http://www.openproblemgarden.org/op/do_any_three_longest_paths_in_a_connected_graph_have_a_vertex_in_common}.
\end{itemize}
\end{document}
